\documentclass[11pt]{article}

\usepackage[a4paper,margin=0.78in]{geometry}
\usepackage[T1]{fontenc}
\usepackage[utf8]{inputenc}
\usepackage{lmodern}
\usepackage{amsmath,amssymb,amsthm,mathtools}
\usepackage{enumitem,booktabs,array,tabularx,microtype,xcolor}
\usepackage[numbers,sort&compress]{natbib}
\usepackage[colorlinks=true,linkcolor=blue!65!black,
  citecolor=blue!65!black,urlcolor=blue!70!black]{hyperref}
\usepackage[nameinlink,noabbrev]{cleveref}
\setlength{\emergencystretch}{3em}

\hypersetup{
  pdftitle={MVP7 on the Occurrence--Phase--Atomic Corridor},
  pdfauthor={Liang Wang},
  pdfsubject={A dynamic source-locked route decision after metric all-prefix return},
  pdfkeywords={MVP7, route classifier, occurrence crosswalk, phase metric, fixed atom}
}

\newtheorem{theorem}{Theorem}[section]
\newtheorem{proposition}[theorem]{Proposition}
\newtheorem{corollary}[theorem]{Corollary}
\theoremstyle{definition}
\newtheorem{definition}[theorem]{Definition}
\theoremstyle{remark}
\newtheorem{remark}[theorem]{Remark}

\newcommand{\GO}{\textnormal{\textsc{go}}}
\newcommand{\AINF}{\textnormal{\textsc{architecture-infeasible}}}
\newcommand{\REROUTE}{\textnormal{\textsc{reroute}}}
\newcommand{\STOPR}{\textnormal{\textsc{stop-route}}}
\newcommand{\NT}{\textnormal{\textsc{not-testable}}}
\newcommand{\AF}{\textnormal{\textsc{arithmetic-frontier}}}
\newcommand{\OPEN}{\textnormal{\textsc{open}}}
\newcommand{\INVALID}{\textnormal{\textsc{invalid}}}
\newcommand{\Lzero}{\mathrm{L0}}
\newcommand{\Lone}{\mathrm{L1}}
\newcommand{\Ltwo}{\mathrm{L2}}
\newcolumntype{Y}{>{\raggedright\arraybackslash}X}
\newcolumntype{P}[1]{>{\raggedright\arraybackslash}p{#1}}

\title{\textbf{MVP7 on the Occurrence--Phase--Atomic Corridor:}\\
\textbf{A Dynamic Source-Locked Route Decision after}\\
\textbf{Metric All-Prefix Return}}
\author{Liang Wang\\
School of Artificial Intelligence and Automation,\\
Huazhong University of Science and Technology, Wuhan 430074, P.R. China\\
\texttt{wangliang.f@gmail.com}}
\date{July 2026}

\begin{document}
\maketitle

\begin{abstract}
TPC-171 integrates the TPC-163--170 batch into a typed source-locked
DAG.  The structural gap is no longer one opaque crosswalk node: it
has three independent roots for a theorem-backed local occurrence
family, actual active support, and canonical minimal representation.
Four independent physical registries remain missing.  In parallel,
the determinant-two core now has a fixed-\(X\)-power all-prefix
packet theorem for Lebesgue-almost every fixed phase.  That theorem
does not reach a named physical phase.

MVP7 turns this state into a dynamic classifier.  Complete
architecture routes and arithmetic method subroutes inhabit disjoint
registries.  Only a stopped architecture can trigger rerouting, and
only through a fresh registry and a source-backed typed crosswalk.
The uncontrolled metric-to-atomic nonimplication stops one arithmetic
method cell; it does not stop the direct fixed-phase twist or the
selected occurrence architecture.  A six-axis projection records all
remaining quantifier changes.

After source, schema, type, route, frontier and Endpoint V4 validation,
the current verdict is
\[
 \mathsf{NOT\_TESTABLE}.
\]
The selected first missing object is the source-backed local
occurrence-edge family, while the full seven-node blocker antichain
is preserved.  The arithmetic advance remains
\(\Lone\) actual-core phase metric, not program-positive \(\Ltwo\).
No strict \(1/400\), prime-pair lower bound or twin-prime theorem is
claimed.
\end{abstract}

\section{Imported state}

Let \(\mathcal M_{171}\) be the canonical TPC-171 manifest.  MVP7
locks its generator, manifest, audit, schemas and paper source, and
also locks the two exact scoped-stop sources used by the classifier.
Every hash has integrity semantics only.

The selected map root has the minimal not-testable antichain
\[
\begin{split}
\mathcal B_{\rm NT}=\{&
\mathsf{H1.source\_backed\_local\_occurrence\_edge\_family},\\
&\mathsf{H1.actual\_active\_support\_certificate},\\
&\mathsf{H1.canonical\_minimal\_representation\_certificate},\\
&\mathsf{H9.literal\_weight\_registry},
\mathsf{H9.phase\_cell\_registry},\\
&\mathsf{H9.endpoint\_registry},
\mathsf{H9.normalization\_registry}\}.
\end{split}
\tag{1}\label{eq:blockers}
\]
The first element is a reporting representative selected by the
declared obligation order.  It does not replace the other six.

The separate parent-ready frontier is
\[
\mathcal F_{\rm open}
=\{\mathsf{O161.bad\_endpoint\_pointwise\_fixed\_atom},
\mathsf{O161.direct\_additive\_twist\_fixed\_atom}\}.
\tag{2}\label{eq:open}
\]
The sets in \eqref{eq:blockers} and \eqref{eq:open} are disjoint in
both status and role.  A missing artifact and an open theorem target
are not interchangeable.

On the arithmetic side the strongest imported result is
\[
 \begin{gathered}
 \mathsf{A170.metric\_packet\_corridor}\\
 =
 \mathsf{PROVED\_L1\_ACTUAL\_CORE\_PHASE\_METRIC\_}\\[-2pt]
 \mathsf{PACKET\_CORRIDOR}.
 \end{gathered}
\tag{3}\label{eq:A170}
\]
It is maximal over every prefix in each declared terminal shell and,
on the prescribed packet schedule, holds for Lebesgue-almost every
fixed phase with every exponent \(\delta<1/4\).  Its analytic norm is
\(L^2_{\rm phase,maximal,BC}\).  The program-positive \(\Ltwo\) flag,
named-fixed-atom flag and production-phase-registry flag are all
false.

\section{Six-axis projection}

For each theorem the TPC-171 contract records the signature
\[
 \mathfrak q=(\mathfrak c,\mathfrak p,\mathfrak e,
              \mathfrak s,\mathfrak d,\mathfrak a),
\tag{4}\label{eq:q}
\]
with carrier, phase, endpoint, scale, decay and support coordinates.
MVP7 displays the strongest proved signature and the physical
signature required at the endpoint:
\[
\begin{array}{c|c|c}
\text{axis}&\text{proved}&\text{required}\\
\hline
\mathfrak c&\text{explicit packet corridor}&
             \text{actual fixed-}h_0\text{ packet}\\
\mathfrak p&\text{Lebesgue-a.e.\ fixed phase}&
             \text{named fixed atom}\\
\mathfrak e&\text{all prefixes in a }\theta\text{-shell}&
             \text{deterministic all-prefix}\\
\mathfrak s&\text{eventually on prescribed schedule}&
             \text{deterministic all-scale}\\
\mathfrak d&\text{fixed-}X\text{ phase-metric power}&
             \text{fixed-}X\text{ fixed-atom power}\\
\mathfrak a&\text{actual arithmetic core}&
             \text{actual active support}.
\end{array}
\tag{5}\label{eq:projection}
\]
No coordinate in \eqref{eq:projection} has yet been bridged.  In
particular, a fixed-\(X\) exponent in the fifth row cannot pay for
missing phase or support quantifiers in the second and sixth rows.

\begin{proposition}[Atomic firewall]
\label{prop:atomic}
The TPC-170 nonimplication stops only the method
\[
 \mathsf{phase\_metric\_uncontrolled\_atomic}.
\tag{6}\label{eq:atomic-stop}
\]
It does not stop a source-backed metric-to-fixed-atom bridge, a
pointwise direct-twist theorem, a pointwise bad-endpoint theorem or an
architecture route.
\end{proposition}

\begin{proof}
The nonimplication uses an arbitrary atomic measure concentrated on
the schedule-dependent metric null set.  Its exact scope is
uncontrolled atomic promotion.  A source-backed bridge would add new
information about the physical phase, and a pointwise theorem would
replace the metric quantifier.  Neither lies in the stopped cell.
Architecture routes concern the construction and return of the
physical carrier and have a different route type.
\end{proof}

\section{Two route families}

The classifier contains two registries.

\paragraph{Architecture routes.}
These are complete candidate paths to the physical endpoint:
\[
\mathcal R_{\rm arch}
=\{\mathsf{current\ archive},\mathsf{occurrence\ map+},
    \mathsf{scalar+ETO}\}.
\tag{7}\label{eq:arch}
\]
The current-archive-only cell is stopped in the exact scope
\(\mathsf{CURRENT\_ARCHIVE\_FIELDS\_ONLY}\).  The selected
occurrence-augmented map and alternative scalar-plus-ETO architectures
remain \(\NT\).  Their universe is not proved complete.

\paragraph{Arithmetic method subroutes.}
These are theorem choices inside an architecture:
\[
\begin{gathered}
\mathsf{periodic\ major},\quad
\mathsf{direct\ twist\ fixed\ atom},\quad
\mathsf{phase\ metric\ source\ backed},\\
\mathsf{phase\ metric\ uncontrolled\ atomic},\quad
\mathsf{bad\ endpoint\ pointwise\ fixed\ atom}.
\end{gathered}
\tag{8}\label{eq:methods}
\]
No element of \eqref{eq:methods} is architecture-reroute eligible.

\begin{definition}[Legal architecture reroute]
\label{def:reroute}
A reroute is legal only if:
\begin{enumerate}[label=(\roman*),leftmargin=2.2em]
\item the selected complete architecture has an exact source-backed
  stop in matching scope;
\item the alternative is another architecture route, not an
  arithmetic method;
\item the alternative has a genuinely fresh registry; and
\item a source-backed typed crosswalk connects the old and new
  carriers, scopes and normalizations.
\end{enumerate}
\end{definition}

Renaming the stopped registry, citing an arithmetic method or
submitting an unregistered theorem label makes the snapshot
\(\INVALID\), not \(\REROUTE\).

\section{Evidence modes}

The production classifier runs in
\(\mathsf{SOURCE\_LOCKED}\) mode.  Every stop, route-universe theorem
and reroute crosswalk must resolve to a canonical repository path,
hash, claim type, route type and exact scope.  The current production
registry contains:
\begin{enumerate}[leftmargin=2em]
\item the current-archive architecture stop from TPC-154; and
\item the uncontrolled-atomic arithmetic-method stop from TPC-170.
\end{enumerate}
The second record cannot be consumed as architecture evidence.

Classifier reachability is tested in the disjoint
\(\mathsf{SYNTHETIC\_REACHABILITY}\) mode.  Its records carry
\[
 \mathsf{NONE\_SYNTHETIC\_REACHABILITY\_ONLY}
\tag{9}\label{eq:synthetic}
\]
semantics and have no file path or theorem hash.  They prove that the
finite classifier logic reaches each verdict; they prove no
mathematical premise.  Moving one of them into production mode is a
validation failure.

\section{Endpoint V4 classifier state}

Endpoint Ledger V4 retains the six literal gates:
\begin{enumerate}[label=\textnormal{V4.\arabic*},leftmargin=2.8em]
\item literal physical coefficients;
\item fixed physical \(h_0=2\);
\item physical atomic normalization;
\item canonical or minimal representation;
\item actual active support;
\item strict \(1/400\) endpoint slack.
\end{enumerate}
Only V4.2 is currently proved.  The four physical registries and all
other literal gates remain \(\NT\).

Each H9 leaf is a data registry and has decay coordinate
\(\mathsf{NONE}\).  The phase leaf may identify a named phase and the
endpoint leaf may identify deterministic endpoints, but neither fact
is an estimate.  A fixed-\(X\) decay coordinate may appear only on an
independent arithmetic theorem that consumes these records.

The phase-metric charge records
\[
 \delta<\frac14
\quad\text{on a Lebesgue-a.e.\ explicit packet schedule}.
\tag{10}\label{eq:metric-delta}
\]
It is explicitly ineligible for the named-fixed-atom ledger.
Therefore
\[
 \sigma_{\rm named\ atom}=0,\qquad
 \sigma_{\rm required}=\frac1{400}.
\tag{11}\label{eq:sigmas}
\]
The physical loss and strict net slack are undefined.  Unknown
quantities are not zero.

\begin{definition}[Arithmetic-frontier state]
\label{def:AF}
The result \(\AF\) is available only when the complete structural
carrier, production physical registry, named fixed phase and
deterministic physical endpoint are all proved, and every remaining
unresolved node is a scope-matched parent-ready positive-\(\Ltwo\)
arithmetic target.
\end{definition}

Thus the current phase-metric theorem cannot itself trigger
\(\AF\).  It leaves structural and physical artifacts unavailable
and does not have the named-atom quantifier.

\section{Ordered classifier}

Validation occurs before mathematical classification.  Failure of a
source lock, schema, route type, evidence type, frontier invariant or
ledger identity returns the outer result \(\INVALID\).

\begin{definition}[Ordered valid-snapshot verdicts]
\label{def:order}
For a valid snapshot, the first true predicate in
\[
\boxed{
\begin{gathered}
\GO,\quad\AINF,\quad\REROUTE,\quad\STOPR,\\
\NT,\quad\AF,\quad\OPEN
\end{gathered}
}
\tag{12}\label{eq:order}
\]
is returned.
\begin{enumerate}[label=(\roman*),leftmargin=2.2em]
\item \(\GO\): every active requirement is proved in matching scope,
  all literal gates pass, and the exact endpoint slack is strictly
  positive.
\item \(\AINF\): an independent source-backed theorem proves the
  architecture universe complete and every complete architecture
  cell has an exact stop.
\item \(\REROUTE\): the selected architecture is stopped and a legal
  alternative satisfies \cref{def:reroute}.
\item \(\STOPR\): the selected architecture is stopped and no legal
  alternative is selected.
\item \(\NT\): the selected-root typed blocker antichain is nonempty.
\item \(\AF\): the conditions of \cref{def:AF} hold.
\item \(\OPEN\): every other valid state.
\end{enumerate}
\end{definition}

\begin{theorem}[Classifier totality and reachability]
\label{thm:total}
Validation followed by \eqref{eq:order} returns exactly one of
\[
\begin{gathered}
\INVALID,\quad\GO,\quad\AINF,\quad\REROUTE,\\
\STOPR,\quad\NT,\quad\AF,\quad\OPEN.
\end{gathered}
\]
Every result is reached by a finite typed regression state.
\end{theorem}

\begin{proof}
Validation is deterministic and either fails or produces a valid
state.  On a valid state the ordered list has a final catch-all, so a
unique first true predicate exists.  The synthetic suite supplies one
separate assumed-predicate state for each valid verdict and one
malformed route-type state for \(\INVALID\).  The evidence-mode
firewall \eqref{eq:synthetic} prevents these reachability witnesses
from entering the production result.
\end{proof}

\section{MVP7 decision}

\begin{theorem}[Source-locked MVP7 verdict]
\label{thm:mvp7}
For the TPC-163--171 source-locked state,
\[
\boxed{
\begin{gathered}
\operatorname{Verdict}_{163:171}=\NT,\\
\mathsf{FirstMissing}
=\mathsf{H1.source\_backed\_local\_}\\[-2pt]
\mathsf{occurrence\_edge\_family}.
\end{gathered}
}
\tag{13}\label{eq:mvp7}
\]
The complete blocker antichain is \eqref{eq:blockers}; the two
pointwise arithmetic targets \eqref{eq:open} remain parent-ready
\(\OPEN\).
\end{theorem}

\begin{proof}
All production source locks, route types, evidence types, frontier
lists and Endpoint V4 identities validate, excluding \(\INVALID\).
The selected occurrence map architecture has no source-backed stop,
and architecture-universe completeness is unproved.  Hence neither
\(\AINF\), \(\REROUTE\) nor \(\STOPR\) applies.  The nonempty
selected-root antichain \eqref{eq:blockers} then triggers \(\NT\)
before \(\AF\) or \(\OPEN\).

Independently, the phase projection \eqref{eq:projection} is
incomplete on every axis, the named-atom exponent in
\eqref{eq:sigmas} is zero and five literal gates remain unknown.
Thus \(\GO\) is also false.  The declared obligation order chooses
the first element of \eqref{eq:blockers} without discarding the other
minimal blockers.
\end{proof}

\begin{corollary}[Exact progress label]
The current state has:
\begin{enumerate}[leftmargin=2em]
\item \(\Lone\) structural localization into three production roots;
\item \(\Lone\) actual-core phase-metric all-prefix packet control;
\item no actual active-support certificate;
\item no named-fixed-phase theorem; and
\item no program-positive \(\Ltwo\) or strict \(1/400\) endpoint.
\end{enumerate}
\end{corollary}

\section{Mutation audit and next decisions}

The audit regenerates the snapshot dynamically from TPC-171.  It
rejects source-hash drift and proof semantics; a Lebesgue-a.e.\ phase
promoted to a named atom; a metric exponent promoted to
program-\(\Ltwo\); a metric charge entered as fixed-atom power; an
arithmetic method inserted into the architecture universe; an
arithmetic method marked reroute eligible; route-universe completeness
without evidence; a fabricated stop label; the uncontrolled-atomic
stop promoted to a direct-twist stop; a collapsed blocker antichain;
an \(\OPEN/\NT\) frontier merge; a positive named-atom exponent
without a bridge; and a \(1/400\) pass without the literal gates.

The next structural construction must supply a theorem-backed local
occurrence-edge family, with source locators and compatible overlaps.
Actual active support and canonical minimal representation remain
parallel obligations, not consequences.  The arithmetic branch has
two live choices: prove the production named phase lies in the
TPC-170 good set through a source-backed bridge, or establish a
pointwise fixed-phase theorem.  If one choice fails in its exact
method cell, the other remains open.

\section{Conclusion}

MVP7 confirms that the route is still live but not yet testable on the
physical carrier.  The new arithmetic corridor is real and stronger
than the earlier almost-endpoint theorem: it reaches all prefixes
with a fixed-\(X\) metric exponent along an explicit schedule.  The
remaining gap is now a sharp quantifier interface, not an unspecified
``phase problem.''  Until a named physical atom and the three
structural roots are supplied, the only source-compatible verdict is
\(\NT\).  This statement neither implies nor suggests a prime-pair
lower bound or the twin-prime conjecture.

\nocite{*}
\bibliographystyle{plainnat}
\bibliography{references}

\end{document}
